\section*{Transport beyond the free field}
The free Dirac construction resolves a translation-invariant field into independent Fourier modes. Gauge coupling changes that organization. Transport between neighboring fibers carries additional information, and a comparison between resolutions must preserve the order and orientation of the links.

A link has a named source and target. Its transformation law is fixed by that direction. Composition cancels the gauge factors at internal endpoints, while a closed path retains conjugation at its basepoint. This elementary mechanism supports the treatment of holonomy, curvature, covariant differences, and refinement.

The finite Maxwell--Dirac example combines matter continuity, electric work, and propagation of the Gauss constraint. Geometric interpolation then separates edge compatibility, curvature defects, and Hilbert normalization. Regularity assumptions turn these structural conditions into quantitative estimates.

The constructions concern specified finite systems and prescribed background connections. A coupled continuum gauge theory additionally requires an operator domain, existence of the field evolution, propagation of constraints, and a convergence theorem in an explicitly chosen topology. Local covariance does not discharge those obligations.

The visual atlas develops the finite examples. The technical appendix records formal source coverage and the limits of its verification evidence. The argument here remains readable without following implementation identifiers.

\clearpage
\section*{Reading routes and conventions}
\addcontentsline{toc}{section}{Reading routes and conventions}
Begin with endpoint covariance if gauge theory is new to you. The cancellation of intermediate frame changes explains both open-path covariance and closed-loop conjugation. In an Abelian group, conjugation disappears; in a non-Abelian group, the trace or conjugacy class is the natural loop observable. The ordered product always records which link acts first.

For finite dynamics, continue to the semidiscrete matter--electric-field system. Here continuity, preservation of a Gauss residual, and energy exchange follow from a specified finite system of equations. The distinction between a time-dependent frame change and the residual symmetry after choosing temporal gauge is made explicit. The consistency calculation also states how a smooth connection is converted into a link; its sampling rule affects the order of accuracy.

For refinement, follow the second part from actual half-edge geometry to cubical paths and Hilbert norms. The text keeps the straight-edge mismatch separate from the curvature contribution. Regular line-integral links remove the former and bound the latter. The reconstruction discussion then explains what is known about state spaces and which generator or evolution comparisons still have to be proved.

\begin{longtable}{@{}L{.24\linewidth}L{.67\linewidth}@{}}
\toprule
Convention & Meaning\\\midrule\endhead
Inner product & $\inner{v}{w}$ is conjugate-linear in $v$ and linear in $w$. Matrix adjoints are written $T^*$.\\
Endpoint transport & $T_{y\leftarrow x}$ maps the fiber at $x$ to the fiber at $y$. Under a local unitary frame $g$, it becomes $g_yT_{y\leftarrow x}g_x^{-1}$.\\
Matrix order & $T_{z\leftarrow y}T_{y\leftarrow x}$ transports first from $x$ to $y$, then from $y$ to $z$. The rightmost factor acts first.\\
Linked shifts & A coefficient $U_j$ in $U_j\psi_{j+1}$ transports from $j+1$ back to $j$. Thus $U_j'=g_jU_jg_{j+1}^{-1}$. An outward link uses the opposite endpoint convention.\\
Gauge groups & $U(1)$ links are unit complex numbers. The finite non-Abelian construction uses unitary matrices, with $SU(2)$ as an explicit example. Scalar multiplication identities are not silently generalized to noncommuting matrices.\\
Mesh and time & Mesh widths are positive when dividing by them. The finite matter system uses $\hbar=1$ and the displayed charge convention. Refinement level and evolution time are different parameters.\\
Norm and measure & Spinor and operator norms are specified in the statement. Half-edge and cubical interpolation use their stated cell weights; an unweighted duplication map is not automatically isometric.\\
Background field & A prescribed connection supplies link integrals and curvature estimates. A dynamical gauge field additionally needs its own evolution and constraints.\\
\bottomrule
\end{longtable}

The geometric source modules also use forward-link notation. The source notes translate it explicitly into the endpoint convention above. In particular, the word ``first'' in an edge-subdivision structure is read together with its product equation and endpoint convention.

\clearpage
